% Options for packages loaded elsewhere
% Options for packages loaded elsewhere
\PassOptionsToPackage{unicode}{hyperref}
\PassOptionsToPackage{hyphens}{url}
\PassOptionsToPackage{dvipsnames,svgnames,x11names}{xcolor}
%
\documentclass[
  letterpaper,
  DIV=11,
  numbers=noendperiod]{scrartcl}
\usepackage{xcolor}
\usepackage{amsmath,amssymb}
\setcounter{secnumdepth}{-\maxdimen} % remove section numbering
\usepackage{iftex}
\ifPDFTeX
  \usepackage[T1]{fontenc}
  \usepackage[utf8]{inputenc}
  \usepackage{textcomp} % provide euro and other symbols
\else % if luatex or xetex
  \usepackage{unicode-math} % this also loads fontspec
  \defaultfontfeatures{Scale=MatchLowercase}
  \defaultfontfeatures[\rmfamily]{Ligatures=TeX,Scale=1}
\fi
\usepackage{lmodern}
\ifPDFTeX\else
  % xetex/luatex font selection
\fi
% Use upquote if available, for straight quotes in verbatim environments
\IfFileExists{upquote.sty}{\usepackage{upquote}}{}
\IfFileExists{microtype.sty}{% use microtype if available
  \usepackage[]{microtype}
  \UseMicrotypeSet[protrusion]{basicmath} % disable protrusion for tt fonts
}{}
\makeatletter
\@ifundefined{KOMAClassName}{% if non-KOMA class
  \IfFileExists{parskip.sty}{%
    \usepackage{parskip}
  }{% else
    \setlength{\parindent}{0pt}
    \setlength{\parskip}{6pt plus 2pt minus 1pt}}
}{% if KOMA class
  \KOMAoptions{parskip=half}}
\makeatother
% Make \paragraph and \subparagraph free-standing
\makeatletter
\ifx\paragraph\undefined\else
  \let\oldparagraph\paragraph
  \renewcommand{\paragraph}{
    \@ifstar
      \xxxParagraphStar
      \xxxParagraphNoStar
  }
  \newcommand{\xxxParagraphStar}[1]{\oldparagraph*{#1}\mbox{}}
  \newcommand{\xxxParagraphNoStar}[1]{\oldparagraph{#1}\mbox{}}
\fi
\ifx\subparagraph\undefined\else
  \let\oldsubparagraph\subparagraph
  \renewcommand{\subparagraph}{
    \@ifstar
      \xxxSubParagraphStar
      \xxxSubParagraphNoStar
  }
  \newcommand{\xxxSubParagraphStar}[1]{\oldsubparagraph*{#1}\mbox{}}
  \newcommand{\xxxSubParagraphNoStar}[1]{\oldsubparagraph{#1}\mbox{}}
\fi
\makeatother


\usepackage{longtable,booktabs,array}
\newcounter{none} % for unnumbered tables
\usepackage{calc} % for calculating minipage widths
% Correct order of tables after \paragraph or \subparagraph
\usepackage{etoolbox}
\makeatletter
\patchcmd\longtable{\par}{\if@noskipsec\mbox{}\fi\par}{}{}
\makeatother
% Allow footnotes in longtable head/foot
\IfFileExists{footnotehyper.sty}{\usepackage{footnotehyper}}{\usepackage{footnote}}
\makesavenoteenv{longtable}
\usepackage{graphicx}
\makeatletter
\newsavebox\pandoc@box
\newcommand*\pandocbounded[1]{% scales image to fit in text height/width
  \sbox\pandoc@box{#1}%
  \Gscale@div\@tempa{\textheight}{\dimexpr\ht\pandoc@box+\dp\pandoc@box\relax}%
  \Gscale@div\@tempb{\linewidth}{\wd\pandoc@box}%
  \ifdim\@tempb\p@<\@tempa\p@\let\@tempa\@tempb\fi% select the smaller of both
  \ifdim\@tempa\p@<\p@\scalebox{\@tempa}{\usebox\pandoc@box}%
  \else\usebox{\pandoc@box}%
  \fi%
}
% Set default figure placement to htbp
\def\fps@figure{htbp}
\makeatother





\setlength{\emergencystretch}{3em} % prevent overfull lines

\providecommand{\tightlist}{%
  \setlength{\itemsep}{0pt}\setlength{\parskip}{0pt}}



 


\KOMAoption{captions}{tableheading}
\usepackage{longtable}
\usepackage{multirow}
\usepackage{array}
\usepackage{booktabs}
\usepackage[table]{xcolor}
\usepackage{hyperref}
\usepackage{fontawesome5}
\usepackage{tikz}
\usetikzlibrary{arrows.meta,positioning}
\makeatletter
\@ifpackageloaded{caption}{}{\usepackage{caption}}
\AtBeginDocument{%
\ifdefined\contentsname
  \renewcommand*\contentsname{Table of contents}
\else
  \newcommand\contentsname{Table of contents}
\fi
\ifdefined\listfigurename
  \renewcommand*\listfigurename{List of Figures}
\else
  \newcommand\listfigurename{List of Figures}
\fi
\ifdefined\listtablename
  \renewcommand*\listtablename{List of Tables}
\else
  \newcommand\listtablename{List of Tables}
\fi
\ifdefined\figurename
  \renewcommand*\figurename{Figure}
\else
  \newcommand\figurename{Figure}
\fi
\ifdefined\tablename
  \renewcommand*\tablename{Table}
\else
  \newcommand\tablename{Table}
\fi
}
\@ifpackageloaded{float}{}{\usepackage{float}}
\floatstyle{ruled}
\@ifundefined{c@chapter}{\newfloat{codelisting}{h}{lop}}{\newfloat{codelisting}{h}{lop}[chapter]}
\floatname{codelisting}{Listing}
\newcommand*\listoflistings{\listof{codelisting}{List of Listings}}
\makeatother
\makeatletter
\makeatother
\makeatletter
\@ifpackageloaded{caption}{}{\usepackage{caption}}
\@ifpackageloaded{subcaption}{}{\usepackage{subcaption}}
\makeatother
\usepackage{bookmark}
\IfFileExists{xurl.sty}{\usepackage{xurl}}{} % add URL line breaks if available
\urlstyle{same}
\hypersetup{
  colorlinks=true,
  linkcolor={blue},
  filecolor={Maroon},
  citecolor={Blue},
  urlcolor={Blue},
  pdfcreator={LaTeX via pandoc}}


\author{}
\date{}
\begin{document}



\begin{longtable*}{p{3cm} p{11.5cm}} \caption[Principales hitos en el EEES]{Principales hitos en la instauración y desarrollo del Espacio Europeo de Educación Superior (EEES). Fuente: \protect\url{https://ehea.info/}.\label{tbl-EEES}}\\

\rowcolors{2}{lightgray}{white}
\toprule
\textbf{Año y lugar} & \textbf{Hitos principales} \\
\midrule
\endfirsthead

\hline
\textbf{Año y lugar} & \textbf{Hitos principales} \\
\hline
\endhead

\hline
\endfoot

\hline
\multicolumn{2}{r}{\textit{Continúa en la página siguiente}} \\
\endfoot

\hline
\endlastfoot

1998, París &
\begin{itemize}
\item \textbf{Declaración de la Sorbona.} Francia, Alemania, Italia y Reino Unido impulsan la armonización de la arquitectura del sistema europeo de educación superior.
\item Se sientan las bases para la adopción de un sistema estructurado en ciclos y para el desarrollo posterior del EEES.
\end{itemize}
\\ \hline

1999, Bolonia &
\begin{itemize}
\item \textbf{Declaración de Bolonia.} 29 países europeos firman el acuerdo fundacional del Proceso de Bolonia.
\item Se establecen los objetivos de crear el EEES antes de 2010, promover la movilidad, el reconocimiento de titulaciones y la comparabilidad de los sistemas.
\end{itemize}
\\ \hline

2001, Praga &
\begin{itemize}
\item Los ministros de Educación revisan el proceso iniciado en Bolonia y fijan directrices y prioridades para los años siguientes, reafirmando el compromiso de constituir el EEES en 2010.
\end{itemize}
\\ \hline

2003, Graz y Berlín &
\begin{itemize}
\item En la II Convención de Instituciones de Educación Superior y en la conferencia ministerial de Berlín se refuerza la dimensión institucional del proceso y su vinculación con el Espacio Europeo de Investigación.
\item Se consolidan además diversas estructuras de seguimiento del proceso.
\end{itemize}
\\ \hline

2005, Bergen &
\begin{itemize}
\item Se subraya la importancia de implicar activamente a estudiantes, instituciones de educación superior, profesorado y empleadores.
\item Se presta especial atención al desarrollo de la investigación y de los programas de doctorado como tercer ciclo.
\end{itemize}
\\ \hline

2007, Londres &
\begin{itemize}
\item El proceso alcanza 46 países participantes.
\item Se evalúan los progresos realizados en movilidad, estructura de titulaciones, reconocimiento, marcos de cualificaciones, aprendizaje a lo largo de la vida, aseguramiento de la calidad y dimensión social.
\item Se fijan como prioridades futuras la movilidad, la dimensión social, la recogida de datos y la empleabilidad.
\end{itemize}
\\ \hline

2009, Lovaina/ Louvain-la-Neuve &
\begin{itemize}
\item Se establecen las áreas principales de trabajo para la década siguiente.
\item El proceso entra en una fase más centrada en la implantación efectiva de las reformas y en su profundización.
\end{itemize}
\\ \hline

2010, Budapest y Viena &
\begin{itemize}
\item Se celebra la conferencia aniversario de los diez años del Proceso de Bolonia.
\item Se produce el lanzamiento oficial del EEES.
\end{itemize}
\\ \hline

2012, Bucarest &
\begin{itemize}
\item El Comunicado de Bucarest se centra en consolidar el EEES en un contexto de crisis económica.
\item Se fijan como prioridades la educación superior de calidad para todos, la mejora de la empleabilidad y la movilidad.
\end{itemize}
\\ \hline

2015, Ereván &
\begin{itemize}
\item Se refuerza el seguimiento de los compromisos esenciales del proceso.
\item Se aprueba la revisión de los ESG 2015.
\end{itemize}
\\ \hline

2018, París &
\begin{itemize}
\item Se insiste en la aplicación efectiva de los compromisos clave de Bolonia.
\item Se subrayan valores fundamentales, inclusión, empleabilidad, movilidad y calidad.
\end{itemize}
\\ \hline

2020, Roma &
\begin{itemize}
\item Visión de un EEES más inclusivo, innovador e interconectado hacia 2030.
\item Anexos sobre libertad académica, dimensión social y mejora del aprendizaje.
\end{itemize}
\\ \hline

2024, Tirana &
\begin{itemize}
\item Evaluación de avances en el 25.º aniversario del Proceso de Bolonia.
\item Reafirmación de pilares: tres ciclos, ECTS, reconocimiento y cultura de calidad.
\end{itemize}
\\ \hline

\end{longtable*}




\end{document}
